\documentclass[../DoAn.tex]{subfiles}
\begin{document}

\begin{center}
    \Large{\textbf{TÓM TẮT NỘI DUNG ĐỒ ÁN}}\\
\end{center}
\vspace{1cm}

Trong bối cảnh chuyển đổi số đang diễn ra mạnh mẽ trong lĩnh vực giáo dục, các trung tâm giảng dạy ngoại ngữ ngày càng chú trọng đến việc ứng dụng công nghệ thông tin nhằm nâng cao hiệu quả quản lý và vận hành. Bên cạnh hoạt động đào tạo chuyên môn, công tác quản lý tài sản đóng vai trò quan trọng trong việc đảm bảo sử dụng hiệu quả cơ sở vật chất, tiết kiệm chi phí và hỗ trợ ra quyết định cho nhà quản lý. Tuy nhiên, trên thực tế, nhiều trung tâm ngoại ngữ vẫn đang quản lý tài sản bằng các phương pháp thủ công hoặc các công cụ đơn giản như bảng tính, dẫn đến nhiều hạn chế như dữ liệu thiếu đồng bộ, khó theo dõi vòng đời tài sản, không kiểm soát được lịch sử sử dụng và bảo trì, cũng như gặp khó khăn trong công tác thống kê và báo cáo.

Trước thực trạng đó, đồ án lựa chọn hướng tiếp cận xây dựng một hệ thống quản lý tài sản. Việc lựa chọn hướng tiếp cận này nhằm tận dụng các chức năng quản lý tài sản sẵn có, giảm chi phí triển khai, đồng thời cho phép tùy biến linh hoạt để phù hợp với đặc thù của trung tâm giảng dạy ngoại ngữ. Trên cơ sở nền tảng được lựa chọn, đồ án tiến hành phân tích yêu cầu thực tế, thiết kế mô hình quản lý tài sản và triển khai hệ thống với các chức năng chính như quản lý danh mục tài sản, phân loại tài sản theo phòng học và mục đích sử dụng, theo dõi tình trạng, lịch sử sử dụng và hỗ trợ phân quyền người dùng.

Đóng góp chính của đồ án là việc xây dựng và triển khai thử nghiệm thành công một hệ thống quản lý tài sản có tính ứng dụng thực tiễn cao cho trung tâm giảng dạy ngoại ngữ. Kết quả đạt được cho thấy hệ thống giúp chuẩn hóa quy trình quản lý tài sản, nâng cao độ chính xác của dữ liệu, hỗ trợ hiệu quả cho công tác thống kê, báo cáo và góp phần cải thiện hiệu quả quản lý và vận hành của trung tâm.

\begin{flushright}
Sinh viên thực hiện\\
\begin{tabular}{@{}c@{}}
\textit{(Ký và ghi rõ họ tên)}
\end{tabular}
\end{flushright}

\end{document}
